\documentclass[11pt]{article}
\usepackage[margin=1.1in]{geometry}
\usepackage{amsmath, amssymb, amsthm}
\usepackage{bm}
\usepackage{hyperref}
\usepackage{booktabs}
\usepackage{parskip}
\usepackage[T1]{fontenc}
\usepackage{lmodern}

\hypersetup{colorlinks=true, linkcolor=blue, urlcolor=blue}

\title{Analytical Fourier Power Spectra of\\
       the Lorentzian and Logistic Dielectric Profiles}
\author{Spatial Kramers--Kronig Project}
\date{}

\newcommand{\kk}{k_s}
\newcommand{\gam}{\gamma}
\newcommand{\eps}{\varepsilon}
\newcommand{\epsh}{\hat{\varepsilon}}
\DeclareMathOperator{\Res}{Res}

\begin{document}
\maketitle
\tableofcontents
\bigskip

\noindent
These notes derive, from first principles, the analytical forms of the Fourier
power spectra $|\epsh(k)|^2$ for two dielectric profiles that appear in the
spatial Kramers--Kronig (sKK) framework:

\begin{enumerate}
    \item A \textbf{Lorentzian} complex profile $\eps(x) = \eps' + i\eps''$
          analytic in the upper half complex $x$-plane.
    \item A \textbf{logistic} profile whose imaginary part is the Hilbert
          transform of the real part.
\end{enumerate}

\noindent
Throughout, the Fourier transform convention is
\begin{equation}\label{eq:ft_convention}
    \epsh(k) = \int_{-\infty}^{\infty} \eps(x)\, e^{-ikx}\, dx,
    \qquad
    \eps(x) = \frac{1}{2\pi}\int_{-\infty}^{\infty} \epsh(k)\, e^{ikx}\, dk,
\end{equation}
where $k$ is the \emph{angular} spatial frequency ($\mathrm{rad\,\mu m^{-1}}$),
consistent with NumPy's FFT convention \texttt{k = 2*pi*fftfreq(...)}.

% ============================================================
\section{The Lorentzian Complex Profile}
% ============================================================

\subsection{Profile definition}

The Lorentzian dielectric profile is
\begin{equation}\label{eq:lor_profile}
    \eps(x) = n_b^2 - \frac{a\gam}{x + i\gam},
    \qquad a,\gam > 0,
\end{equation}
with real part $\eps'(x)$ and imaginary part $\eps''(x)$ given by
\begin{equation}
    \eps'(x)  = n_b^2 - \frac{a\gam x}{x^2+\gam^2},
    \qquad
    \eps''(x) = \frac{a\gam^2}{x^2+\gam^2}.
\end{equation}
The profile is constructed so that $\eps(z)$ is analytic in the upper
half $z$-plane $(\mathrm{Im}(z)>0)$; the single pole lies at
$z = -i\gam$, strictly in the \emph{lower} half-plane.

\subsection{Fourier transform via contour integration}

Because $n_b^2$ contributes only a $\delta(k)$ term to $\epsh(k)$, we focus on
the non-trivial part
\begin{equation}
    \epsh(k) = -a\gam \int_{-\infty}^{\infty}
        \frac{e^{-ikx}}{x + i\gam}\, dx, \qquad k \neq 0.
\end{equation}
Promote $x$ to a complex variable $z = x + iy$ and close the real-axis
integral with a large semicircular arc of radius $R \to \infty$.

\paragraph{Choosing the contour.}
On the arc $z = Re^{i\theta} = R\cos\theta + iR\sin\theta$, the exponent in
the integrand is
\begin{equation}
    -ikz = -ik(R\cos\theta + iR\sin\theta)
         = \underbrace{kR\sin\theta}_{\mathrm{Re}(-ikz)} \;-\; ikR\cos\theta,
\end{equation}
so $\left|e^{-ikz}\right| = e^{\mathrm{Re}(-ikz)} = e^{k\,\mathrm{Im}(z)} = e^{kR\sin\theta}$.
\begin{itemize}
    \item \textbf{Lower half-plane} (LHP, $\theta\in(-\pi,0)$, $\sin\theta < 0$):
          $|e^{-ikz}| = e^{kR\sin\theta} \to 0$ for $k > 0$.
    \item \textbf{Upper half-plane} (UHP, $\theta\in(0,\pi)$, $\sin\theta > 0$):
          $|e^{-ikz}| = e^{kR\sin\theta} \to 0$ for $k < 0$.
\end{itemize}
Hence the arc integral vanishes by Jordan's lemma in the appropriate half-plane.

\paragraph{Case $k > 0$ (close in the LHP).}
The contour is traversed \emph{clockwise}, enclosing the pole at $z_0 = -i\gam$:
\begin{equation}
    \epsh(k) = -a\gam \cdot (-2\pi i)\,
        \underbrace{\Res_{z=-i\gam}\!\frac{e^{-ikz}}{z+i\gam}}_{\displaystyle= e^{-ik(-i\gam)} = e^{-k\gam}}.
\end{equation}
The residue of $e^{-ikz}/(z+i\gam)$ at $z=-i\gam$ is found by noting that
$(z+i\gam)$ has a simple zero there:
\begin{equation}
    \Res_{z=-i\gam}\frac{e^{-ikz}}{z+i\gam}
    = \lim_{z\to -i\gam}(z+i\gam)\cdot\frac{e^{-ikz}}{z+i\gam}
    = e^{-ik(-i\gam)} = e^{-k\gam}.
\end{equation}
Therefore
\begin{equation}\label{eq:lor_epsh_pos}
    \boxed{\epsh(k) = 2\pi i\, a\gam\, e^{-k\gam}, \qquad k > 0.}
\end{equation}

\paragraph{Case $k < 0$ (close in the UHP).}
The contour is traversed \emph{counter-clockwise}. There are \emph{no} poles
in the UHP (the only pole is $z_0=-i\gam$ in the LHP), so by the residue
theorem
\begin{equation}
    \epsh(k) = 0, \qquad k < 0.
\end{equation}
This confirms that the Lorentzian profile is an analytic signal: all spectral
weight lies at positive spatial frequencies, consistent with $\eps(z)$ being
analytic above the real axis.

\subsection{Power spectrum}

From Eq.~\eqref{eq:lor_epsh_pos},
\begin{equation}\label{eq:lor_power}
    \boxed{|\epsh(k)|^2 = 4\pi^2 a^2 \gam^2\, e^{-2\gam k}, \qquad k > 0.}
\end{equation}
The decay rate $2\gam$ is set by the imaginary distance from the real axis to
the nearest pole. In the numerical scripts this exact formula is implemented as:
\begin{verbatim}
    C_exact = 4 * np.pi**2 * a**2 * gam**2
    pwr_exp_an = C_exact * np.exp(-2 * gam * k_an)
\end{verbatim}
with parameters $a=1$, $\gam=0.01\,\mu\mathrm{m}$, $n_b=1.5$.

% ============================================================
\section{The Logistic Profile}
% ============================================================

\subsection{Profile definition}

The logistic dielectric profile is
\begin{equation}\label{eq:log_profile}
    \eps'(x) = 1 + \frac{\Delta}{1 + e^{\kk x}},
    \qquad \Delta = n_b^2 - 1,\; \kk > 0,
\end{equation}
which interpolates smoothly from $\eps'(-\infty)=1+\Delta=n_b^2$ to
$\eps'(+\infty)=1$ over a characteristic width $1/\kk$.
Unlike the Lorentzian, $\eps'(x)$ is real-valued and is \emph{not} by itself
a valid sKK profile.  The imaginary part $\eps''(x)$ must be added as the
Hilbert transform of $\eps'(x)$.

\subsection{Analytic signal: eliminating \texorpdfstring{$k<0$}{k<0} components}

\paragraph{Hilbert transform in Fourier space.}
For any real-valued $f(x)$, the Hilbert transform $g = \mathcal{H}[f]$ has the
Fourier representation
\begin{equation}
    \hat{g}(k) = -i\,\mathrm{sgn}(k)\,\hat{f}(k).
\end{equation}
Form the \emph{analytic signal} $\eps(x) = \eps'(x) + i\eps''(x)$ where
$\eps'' = \mathcal{H}[\eps']$.  Its Fourier transform is
\begin{equation}
    \epsh(k)
    = \hat{\eps}'(k) + i\hat{\eps}''(k)
    = \hat{\eps}'(k) + i\bigl(-i\,\mathrm{sgn}(k)\bigr)\hat{\eps}'(k)
    = \hat{\eps}'(k)\bigl(1 + \mathrm{sgn}(k)\bigr).
\end{equation}
Evaluating the factor $(1+\mathrm{sgn}(k))$ case by case:
\begin{equation}\label{eq:analytic_signal}
    \epsh(k) =
    \begin{cases}
        2\hat{\eps}'(k), & k > 0, \\
        0,               & k < 0.
    \end{cases}
\end{equation}
\textbf{Key conclusion}: the $k<0$ sector is identically zero regardless of whether
the complex profile $\eps'(z)$ has poles in the UHP or LHP.  The sKK condition
is enforced by the construction $\eps'' = \mathcal{H}[\eps']$, not by analytic
continuation of $\eps'$ alone.  Therefore we need only compute $\hat{\eps}'(k)$
for $k>0$.

\subsection{Fourier transform of \texorpdfstring{$\eps'(x)$}{eps'(x)} for \texorpdfstring{$k>0$}{k>0}}

Again discarding the constant (which contributes only at $k=0$), we need
\begin{equation}
    \hat{\eps}'(k) = \Delta\int_{-\infty}^{\infty}
        \frac{e^{-ikx}}{1 + e^{\kk x}}\, dx, \qquad k > 0.
\end{equation}

\paragraph{Poles of the integrand.}
The denominator $1+e^{\kk z}=0$ when $e^{\kk z} = -1 = e^{i\pi(2m-1)}$
for $m\in\mathbb{Z}$, giving
\begin{equation}\label{eq:log_poles}
    z_m = \frac{i\pi(2m-1)}{\kk}, \qquad m \in \mathbb{Z}.
\end{equation}
All poles lie on the imaginary axis.  For $k>0$ we close in the LHP
($\mathrm{Im}(z)<0$), which contains the poles with $2m-1<0$, i.e.,
$m\leq 0$:
\begin{equation}
    z_m = -\frac{i\pi(2n-1)}{\kk}, \quad n = 1,2,3,\ldots
    \quad (n \equiv 1-m).
\end{equation}
The nearest LHP pole to the real axis is $z_1 = -i\pi/\kk$, at distance
$\pi/\kk$ (not $1/\kk$); the extra factor of $\pi$ arises from solving
$e^{\kk z}=-1 \Rightarrow \kk z = \ln(-1) = i\pi$.

\paragraph{Residues.}
Each pole $z_m$ is simple.  Using the residue formula for $f(z)/g(z)$ at a
simple zero $z_0$ of $g$:
\begin{equation}
    \Res_{z=z_0}\frac{f(z)}{g(z)} = \frac{f(z_0)}{g'(z_0)}.
\end{equation}
With $f(z)=\Delta e^{-ikz}$ and $g(z)=1+e^{\kk z}$, $g'(z)=\kk e^{\kk z}$:
\begin{equation}
    \Res_{z=z_m}\frac{\Delta e^{-ikz}}{1+e^{\kk z}}
    = \frac{\Delta e^{-ikz_m}}{\kk e^{\kk z_m}}
    = \frac{\Delta e^{-ikz_m}}{\kk \cdot (-1)}
    = -\frac{\Delta}{\kk}\,e^{-ikz_m},
\end{equation}
where we used $e^{\kk z_m}=-1$ at the pole.

\paragraph{Summing residues (CW contour).}
For $m \leq 0$ (LHP poles, relabeled as $n=1,2,\ldots$):
\begin{equation}
    \hat{\eps}'(k) = (-2\pi i)
        \sum_{n=1}^{\infty}
        \left(-\frac{\Delta}{\kk}\right)
        e^{-ik\cdot\left(-i\pi(2n-1)/\kk\right)}.
\end{equation}
The exponent becomes
\begin{equation}
    -ik\cdot\left(-\frac{i\pi(2n-1)}{\kk}\right)
    = -\frac{\pi k(2n-1)}{\kk}.
\end{equation}
Therefore
\begin{equation}
    \hat{\eps}'(k) = \frac{2\pi i\,\Delta}{\kk}
        \sum_{n=1}^{\infty} e^{-\pi k(2n-1)/\kk}.
\end{equation}

\paragraph{Evaluating the geometric series.}
Let $\beta \equiv \pi k/\kk > 0$.  Then
\begin{align}
    \sum_{n=1}^{\infty} e^{-\beta(2n-1)}
    &= e^{-\beta} + e^{-3\beta} + e^{-5\beta} + \cdots \notag\\
    &= e^{-\beta}\bigl(1 + e^{-2\beta} + e^{-4\beta}+\cdots\bigr) \notag\\
    &= \frac{e^{-\beta}}{1 - e^{-2\beta}}
     = \frac{1}{e^{\beta} - e^{-\beta}}
     = \frac{1}{2\sinh\beta}
     = \frac{1}{2\sinh(\pi k/\kk)}.
\end{align}
Hence
\begin{equation}
    \hat{\eps}'(k) = \frac{2\pi i\,\Delta}{\kk}
        \cdot\frac{1}{2\sinh(\pi k/\kk)}
    = \frac{\pi i\,\Delta}{\kk\sinh(\pi k/\kk)}.
\end{equation}
Using Eq.~\eqref{eq:analytic_signal}, the full spectrum for $k>0$ is
\begin{equation}\label{eq:log_epsh}
    \epsh(k) = 2\hat{\eps}'(k) = \frac{2\pi i\,\Delta}{\kk\sinh(\pi k/\kk)}.
\end{equation}

\subsection{Power spectrum}

\begin{equation}\label{eq:log_power}
    \boxed{|\epsh(k)|^2 = \frac{C}{\sinh^2(\pi k/\kk)},
    \qquad
    C = \frac{4\pi^2\Delta^2}{\kk^2},
    \qquad k>0,}
\end{equation}
with no free fitting parameters.  In the numerical scripts this exact formula
appears as:
\begin{verbatim}
    Delta = nb_g**2 - 1
    C_exact = 4 * np.pi**2 * Delta**2 / k_steep**2
    pwr_sinh = C_exact / np.sinh(np.pi * k_an / k_steep)**2
\end{verbatim}
with $\kk = k_{\text{steep}} = 100\,\mu\mathrm{m}^{-1}$, $n_b=1.7$.

\subsection{Asymptotic regimes}

The two limiting behaviors of $\sinh(\pi k/\kk)$ yield distinct power-law
and exponential regimes.

\paragraph{Small $k$ (algebraic regime, $k \ll \kk/\pi$).}
Use the Taylor series $\sinh(x)\approx x$ for $x\ll 1$:
\begin{equation}
    \sinh\!\left(\frac{\pi k}{\kk}\right) \approx \frac{\pi k}{\kk},
    \qquad k \ll \frac{\kk}{\pi}.
\end{equation}
Substituting into Eq.~\eqref{eq:log_power}:
\begin{equation}
    |\epsh(k)|^2 \approx
        \frac{4\pi^2\Delta^2/\kk^2}{(\pi k/\kk)^2}
    = \frac{4\Delta^2}{k^2}.
\end{equation}
This $1/k^2$ power law is characteristic of a profile with a finite step
discontinuity (the logistic step of height $\Delta$).

\paragraph{Large $k$ (exponential regime, $k \gg \kk/\pi$).}
For $x\gg 1$, $\sinh(x)\approx e^x/2$:
\begin{equation}
    \sinh\!\left(\frac{\pi k}{\kk}\right) \approx \frac{1}{2}e^{\pi k/\kk},
    \qquad k \gg \frac{\kk}{\pi}.
\end{equation}
Therefore
\begin{equation}
    |\epsh(k)|^2 \approx
        \frac{4\pi^2\Delta^2/\kk^2}{\tfrac{1}{4}e^{2\pi k/\kk}}
    = \frac{16\pi^2\Delta^2}{\kk^2}\,e^{-2\pi k/\kk}
    = 4C\,e^{-2\pi k/\kk}.
\end{equation}
The exponential decay rate is $2\pi/\kk$ — a factor of $\pi$ larger
than the naive estimate $2/\kk$, because the nearest LHP pole of the
logistic is at $\mathrm{Im}(z_1) = -\pi/\kk$, not $-1/\kk$.

\paragraph{Crossover.}
The two asymptotes are equal when
\begin{equation}
    \frac{4\Delta^2}{k^2} = 4C\,e^{-2\pi k/\kk},
\end{equation}
which to leading order gives $k_{\times} \approx \kk/\pi$.
For $\kk=100\,\mu\mathrm{m}^{-1}$ this is $k_{\times}\approx 32\,\mu\mathrm{m}^{-1}$.

% ============================================================
\section{Comparison: Lorentzian vs.\ Logistic}
% ============================================================

\begin{center}
\begin{tabular}{lcc}
\toprule
Property & Lorentzian & Logistic \\
\midrule
Profile & $\eps = n_b^2 - a\gam/(x+i\gam)$ & $\eps' = 1 + \Delta/(1+e^{\kk x})$ \\
Imaginary part & Intrinsic (analytic) & $\eps'' = \mathcal{H}[\eps']$ (imposed) \\
Poles & Single: $z_0 = -i\gam$ & Infinite: $z_n = -i\pi(2n-1)/\kk$ \\
Nearest LHP pole & $\gam$ below real axis & $\pi/\kk$ below real axis \\
Spectrum & $4\pi^2 a^2\gam^2\,e^{-2\gam k}$ & $C/\sinh^2(\pi k/\kk)$ \\
Large-$k$ slope & $-2\gam$ (exp. decay) & $-2\pi/\kk$ \\
Small-$k$ behavior & exponential (no crossover) & $1/k^2$ \\
\bottomrule
\end{tabular}
\end{center}

\noindent
If the transition widths are matched, $\gam = 1/\kk$, the large-$k$ decay
rates differ by $\pi$:
\begin{equation}
    \text{Lorentzian: } -2\gam = -2/\kk,
    \qquad
    \text{Logistic: } -2\pi/\kk.
\end{equation}
The factor of $\pi$ reflects the geometry: the Lorentzian pole sits at
distance $\gam=1/\kk$ from the real axis, whereas the nearest logistic pole
sits at $\pi/\kk$ — a consequence of $\ln(-1)=i\pi$.

\end{document}
